% ID: FIS-OTT-RFL-001
% Titolo: Condizione di riflessione totale
% Disciplina: Fisica
% Area: Ottica
% Argomento: Riflessione
% Sottoargomento: Riflessione totale interna
% Classe: Seconda liceo scientifico
% Difficolta: 3
% Tipo: problema_numerico
% Risultato: n_2 massimo circa 1,29
% Tempo_stimato: 8 min
% Competenze: uso dell'angolo limite, interpretazione fisica, calcolo con indici di rifrazione
% Prerequisiti: legge di Snell, angolo limite, riflessione totale
% Tag: fisica, ottica, riflessione_totale, indice_rifrazione, angolo_limite
% Fonte: verifica_fisica_ottica_anonimizzata
% Autore: Riccardo Carli
% Licenza: CC-BY-SA-4.0

\begin{esercizio}
Un raggio luminoso si propaga in acqua, con indice di rifrazione
\(n_1=1{,}333\), e incide con un angolo di \(75^\circ\) su un secondo mezzo.
Se il raggio viene totalmente riflesso, calcola il massimo valore possibile
dell'indice di rifrazione del secondo mezzo.
\end{esercizio}

\begin{soluzione}
Per avere riflessione totale interna, il raggio deve passare da un mezzo piu
rifrangente a uno meno rifrangente e l'angolo di incidenza deve essere almeno
uguale all'angolo limite.

All'angolo limite vale:
\[
\sin\theta_L=\frac{n_2}{n_1}.
\]
Se la riflessione totale avviene per \(\theta_i=75^\circ\), il massimo valore
possibile di \(n_2\) si ha quando \(\theta_L=75^\circ\):
\[
n_{2,\max}=n_1\sin 75^\circ.
\]

\[
n_{2,\max}=1{,}333\cdot \sin 75^\circ
\simeq 1{,}29.
\]

Il secondo mezzo deve quindi avere indice di rifrazione non superiore a circa
\(1{,}29\).
\end{soluzione}
